\documentclass[11pt]{article}
\usepackage{preamble}
\setlength{\parskip}{4pt}

\begin{document}

\daybanner{AP Statistics \textperiodcentered\ Unit 2 \textperiodcentered\ Topic 2.1 \textperiodcentered\ B: 2026-10-15 \textperiodcentered\ E: 2026-10-28 \textperiodcentered\ TEACHER KEY}{Same Responses, Different Questions}

\sectionbanner{CED TETHER}

{\small
\begin{itemize}[leftmargin=1.5em,itemsep=2pt,topsep=2pt]
  \item \textbf{Skill 2.D} --- Compare distributions or relative positions of points within a distribution.
  \item \textbf{EU UNC-1} --- Graphical representations and statistics allow us to identify and represent key features of data.
  \item \textbf{LO UNC-1.P} --- Compare numerical and graphical representations for two categorical variables. [Skill 2.D]
  \item \textbf{EK UNC-1.P.1} --- Side-by-side bar graphs, segmented bar graphs, and mosaic plots are examples of bar graphs for one categorical variable, broken down by categories of another categorical variable.
  \item \textbf{EK UNC-1.P.2} --- Graphical representations of two categorical variables can be used to compare distributions and/or determine if variables are associated.
  \item \textbf{EK UNC-1.P.3} --- A two-way table, also called a contingency table, is used to summarize two categorical variables. The entries in the cells can be frequency counts or relative frequencies.
  \item \textbf{EK UNC-1.P.4} --- A joint relative frequency is a cell frequency divided by the total for the entire table.
\end{itemize}
}
\textbf{Essential question:} Which display best answers a relationship question when the groups have different sizes?\par
\textbf{DOK-3 skill:} 4.B \textperiodcentered\ \textbf{FRQ pattern:} {\small\texttt{count-\allowbreak{}bars-\allowbreak{}vs-\allowbreak{}segmented-\allowbreak{}bars-\allowbreak{}which-\allowbreak{}answers-\allowbreak{}the-\allowbreak{}question}}\par
\textbf{DOK rationale:} Part (c) weighs competing claims using the day's numerical or graphical evidence and explains a limitation or consequence; the judgment requires linked reasoning beyond a calculation.\par

\frameworkphaseheader{Do Now (first take) $\rightarrow$ Explore (video + follow-along) $\rightarrow$ Exit (finish + turn in)}{1 $\rightarrow$ 3}{45}{%
\begin{itemize}[leftmargin=1.5em,itemsep=2pt,topsep=2pt]
  \item Show the board at the bell and collect a one-sentence first take without confirming a claim.
  \item Run the named video follow-along, then allow about 10 minutes for parts (a)--(c).
  \item Ask for evidence linked to a claim. Collect the sheet and hand-score part (c) with E/P/I.
\end{itemize}
}{%
\begin{itemize}[leftmargin=1.5em,itemsep=2pt,topsep=2pt]
  \item Choose a graph for the relationship question and cite one feature.
  \item Complete the video follow-along about two-way tables and bar-graph displays.
  \item Finish (a)--(c), revise the first take in the exit line, and turn in the sheet.
\end{itemize}
}{%
\begin{itemize}[leftmargin=1.5em,itemsep=2pt,topsep=2pt]
  \item Which number or visible feature supports your claim?
  \item Why does the other claim go beyond that evidence?
  \item What would you tell someone who chose the other claim?
\end{itemize}
}{Circulate and ask students to connect their choice to evidence. Score part (c) using the three required reasoning elements; do not require dropped extension work.}

\begin{callout}{\IconWarn\ WATCH FOR}{calloutred}
\begin{itemize}[leftmargin=1.5em,itemsep=2pt,topsep=2pt]
  \item A choice with copied numbers but no explanation of how they support it.
  \item Treating a model or average as a guaranteed individual outcome.
\end{itemize}
\end{callout}

\sectionbanner{SCORING PART (c) --- E / P / I}

\textbf{Expected elements}
\begin{itemize}[leftmargin=1.5em,itemsep=2pt,topsep=2pt]
  \item \textbf{reasoning-1} (required) --- Chooses B for the relationship using 60 percent versus 30 percent
  \item \textbf{reasoning-2} (required) --- Explains how equal whole-bar heights compare within-grade proportions
  \item \textbf{reasoning-3} (required) --- Acknowledges Andre's group-size point and gives a count question A answers directly
\end{itemize}

{\small\renewcommand{\arraystretch}{1.25}
\noindent\begin{tabular}{|p{0.5in}|p{5.9in}|}
\hline
\textbf{E} & Includes all three required reasoning elements above, with correct contextual evidence. \\ \hline
\textbf{P} & Includes at least one substantive correct reasoning link above, but omits or misstates another required element. \\ \hline
\textbf{I} & Includes no substantive correct reasoning link above; an unsupported choice or copied number alone is insufficient. \\ \hline
\end{tabular}}

\textbf{Common mistakes}
\begin{itemize}[leftmargin=1.5em,itemsep=2pt,topsep=2pt]
  \item Comparing 72 with 18 as if those counts were within-grade percentages
  \item Saying equal-height bars mean the two grade distributions are the same
  \item Claiming Graph A contains no information about percentages rather than saying division is required
  \item Claiming Graph B shows 120 ninth graders and 60 twelfth graders
\end{itemize}

\clearpage
\focusbanner{TODAY'S PROBLEM --- annotated}

Suppose a student council asks 120 ninth graders and 60 twelfth graders whether they prefer Field Day or a Talent Show for the spring celebration. The same responses are shown in two graphs.\par\smallskip \textbf{Andre:} \emph{Graph B hides how many people there are because both bars reach 100\%.}\par \textbf{Leila:} \emph{Graph A hides the relationship because the grade totals are different.}\par
\begin{center}
\begin{tikzpicture}\begin{axis}[width=0.95\linewidth,height=1.05in,ybar stacked,bar width=18pt,ymin=0,ymax=135,ytick={0,60,120},symbolic x coords={9th,12th},xtick=data,enlarge x limits=0.4,title={Graph A},ylabel={Students},tick label style={font=\small},label style={font=\small},title style={font=\small\bfseries},nodes near coords,every node near coord/.append style={font=\small},ymajorgrids]\addplot[fill=calloutblue,draw=dayblue] coordinates {(9th,72) (12th,18)};\addplot[fill=calloutgreen,draw=dayblue] coordinates {(9th,48) (12th,42)};\end{axis}\end{tikzpicture}
\par\smallskip
\begin{tikzpicture}\begin{axis}[width=0.95\linewidth,height=1.05in,ybar stacked,bar width=18pt,ymin=0,ymax=112,ytick={0,50,100},symbolic x coords={9th,12th},xtick=data,enlarge x limits=0.4,title={Graph B},ylabel={Percent},tick label style={font=\small},label style={font=\small},title style={font=\small\bfseries},nodes near coords,every node near coord/.append style={font=\small},ymajorgrids]\addplot[fill=calloutblue,draw=dayblue] coordinates {(9th,60) (12th,30)};\addplot[fill=calloutgreen,draw=dayblue] coordinates {(9th,40) (12th,70)};\end{axis}\end{tikzpicture}
\par{\small Lower: Field Day\quad Upper: Talent Show}
\end{center}
\begin{firsttakebox}{FIRST TAKE (5 min)}
For the question \emph{Is celebration preference related to grade level?}, which graph would you use first? Name one visible feature.\par
\answer{Not graded. The goal is to surface whether students instinctively compare bar totals, segment counts, or within-grade shares.}
\end{firsttakebox}

\textit{Rules callout ``MATCH THE DISPLAY TO THE QUESTION (keep on the page)'' is printed on the student sheet.}\par\medskip

\dokbadge{1}~\textbf{(a)} From Graph A, state the total surveyed in each grade.\par
\answer{There are 120 ninth graders and 60 twelfth graders.}

\dokbadge{2}~\textbf{(b)} Find the percent who prefer Field Day in each grade. Show the fractions and compare the results.\par
\answer{Ninth grade: $72/120=60\%$. Twelfth grade: $18/60=30\%$. Field Day preference is 30 percentage points higher in ninth grade.}

\dokbadge{3}~\textbf{(c)} Which graph better shows whether grade and preference are related? Use both Field Day percentages and explain how the bar heights help. Then explain what Andre is right about and one question Graph A answers more directly. Write 3--4 sentences.\par
\answer{Graph B shows the relationship more directly: its equally tall whole bars let us compare within-grade shares, $60\%$ versus $30\%$, without unequal group sizes distorting the comparison. These different shares suggest an association. Andre is right that Graph B alone does not show how many students were surveyed. Graph A directly answers count questions such as how many twelfth graders chose Talent Show (42).}

\begin{summaryexitbox}{EXIT}
One thing the video changed about my first take: \hrulefill
\end{summaryexitbox}

\end{document}
